\section{Introduction}

Understanding the internal, three-dimensional structure of hadrons, such as the proton, is an important goal in nuclear and particle physics. In this regard, the transverse momentum-dependent (TMD) parton distribution functions (TMDPDFs)~\cite{Collins:1981uk,Collins:1981va} play an important role as they characterize their intrinsic transverse 
partonic structure. These distributions are also essential ingredients in the description of multi-scale and non-inclusive processes, such as Drell-Yan production of electroweak gauge bosons or Higgs bosons or semi-inclusive deep-inelastic scattering with small transverse momentum, in the context of QCD factorization theorems. As a result, they have received considerable attention in the past few decades (for a review, see Ref.~\cite{Angeles-Martinez:2015sea}). More accurate experimental measurements are expected in the coming decades from JLab 12 GeV~\cite{Dudek:2012vr} and the Electron-Ion Colliders in the US~\cite{Accardi:2012qut,AbdulKhalek:2021gbh} and in China~\cite{Anderle:2021wcy}.


In contrast to the TMDPDFs that encode the probability density of parton momenta in hadrons, the transverse momentum-dependent wave functions (TMDWFs) offer a probability amplitude description of the partonic structure of hadrons, from which one can potentially calculate various quark/gluon distributions. In the QCD factorization involving transverse momentum, they are the most important ingredients to predict physical observables in exclusive processes, for instance, weak decays of heavy $B$ meson~\cite{Keum:2000wi,Lu:2000em} which are valuable to extract the CKM matrix element, and to probe new physics beyond the standard model. However, due to the lack of  knowledge of TMDWFs, the one-dimensional lightcone distribution amplitudes (LCDAs) are used instead in most analyses of $B$ decays \cite{Keum:2000wi,Lu:2000em,Nagashima:2002ia}, resulting in uncontrollable errors.  The unprecedented  precision of experimental measurements  of $B$ decays~\cite{Cerri:2018ypt} urgently requires a reliable theoretical  knowledge of TMDWFs. 

A common feature of TMDPDFs and TMDWFs is that they depend both on the longitudinal momentum fraction $x$ and on the transverse spatial separation of partons.  Considerable theoretical efforts have been devoted in recent years to determine these quantities by fitting the pertinent experimental data~\cite{Landry:1999an,Landry:2002ix,DAlesio:2014mrz,Sun:2014dqm,Konychev:2005iy,Bacchetta:2017gcc,Scimemi:2017etj,Scimemi:2019cmh,Bacchetta:2019sam}, which, however, %inevitably introduces systematic uncertainties. Thus 
is limited by the imprecise knowledge of the nonperturbative behaviour of TMDPDFs and TMDWFs. Thus, it is highly desirable to develop a method to calculate them from first-principle approaches such as lattice QCD.

This has been realized in the framework of large momentum effective theory~\cite{Ji:2013dva,Ji:2014gla}, which offers a systematic way to calculate light-cone correlations by simulating time-independent Euclidean correlations on the lattice. Significant progress has been made in calculating various parton quantities from LaMET. For recent reviews, see Ref.~\cite{Cichy:2018mum,Ji:2020ect}. 

A very important result of LaMET development is that the TMDPDFs and TMDWFs can be calculated through the Euclidean quasi-TMDPDFs and quasi-TMDWFs, as well as a universal soft function (factor) ~\cite{Ebert:2019okf,Ji:2019sxk,Ji:2019ewn,Ji:2021znw}. In Ref.~\cite{Ji:2019sxk}, it has been suggested that the form factor of a bi-local four-quark operator calculable on the lattice, can be factorized into quasi-TMDWFs, a universal soft (function) factor and the matching kernel through QCD factorization at large momentum transfer, allowing
for the first time calculating the universal soft function on lattice. Thus the light-cone TMD parton distributions and wave functions can be obtained from numerical calculations of the four-quark form factors and quasi TMDPDFs and TMDWFs on the lattice~\cite{Ji:2019sxk,Ji:2019ewn}.  On the other hand, one can also make use of the QCD factorization to obtain the CS kernel from quasi TMDPDFs and TMDWFs. The first results for the CS kernel based on these proposals have been published recently~\cite{Shanahan:2020zxr,LatticeParton:2020uhz,Schlemmer:2021aij,Li:2021wvl,Shanahan:2021tst}. The quasi TMDWFs approach
to the CS kernel requires two-point function calculations 
and potentially can reach the light-cone limit with relatively
small hadron momenta. 

In this work, we present a state-of-the-art calculation of the CS kernel, based on a lattice QCD analysis of quasi-TMDWFs with $N_f=2+1+1$ valence clover fermions on a staggered quark sea with one-loop matching accuracy. A single ensemble with lattice spacing $a\simeq0.12$~fm,  volume $n_s^3\times n_t=48^3\times64$, and  physical sea-quark masses is used. In order to improve the signal-to-noise ratio, we tune the light-valence quark masses such that $m_\pi = 670$ MeV. The CS  kernel is then extracted through the ratios of the quasi-TMDWFs and the perturbative matching kernels at different momenta, $P^z=2\pi/n_s\times\{8,10,12\}=\{1.72,~2.15,~2.58\}$GeV. This  corresponds to Lorentz boost factors  $\gamma=\{2.57,~3.21,~3.85\}$, respectively. This analysis improves the previous ones~\cite{LatticeParton:2020uhz,Li:2021wvl} by taking into account  the one-loop perturbative contributions, and by analyzing  systematic uncertainties from operator mixing,  higher-order corrections from the scale dependence, and higher power corrections in terms of $1/P^z$.  
 
The remainder of this paper is organized as follows. 
In Sec.~\ref{sec:framework}, we present the theoretical framework to extract the CS kernel from  quasi-TMDWFs. Numerical results for quasi-TMDWFs and CS kernel are presented in Sec.~\ref{sec:numerics}. A brief  summary of this work is given in Sec.~\ref{sec:summary}.  More details about the analysis are collected in the appendix. 
